\section{Вывод термодинамического учёта энтропии}
\label{app:energy}

Выводим Ур.~\eqref{eq:energy-law} и уточняем её точный охват. Принцип
Ландауэра~\cite{landauer1961} гласит, что снижение энтропии системы на $1$~бит
($\ln 2$~нат) при экспорте в термостат температуры $T$ требует не менее
$k_B T\ln 2$ свободной энергии. Логико-термодинамический бюджет тика тем самым
ограничен снизу величиной $k_B T\ln 2$, умноженной на число экспортированных бит
энтропии фон Неймана.

\paragraph{Три члена.} Разложим тик (Ур.~\ref{eq:lomega}), измеряя энтропию фон
Неймана в битах, $S_{\mathrm{vN}}(\Gam)=-\Tr(\Gam\log_2\Gam)$:
\begin{itemize}[leftmargin=1.5em,itemsep=1pt]
  \item \textbf{Унитарный} $-i[H_{\mathrm{eff}},\Gam]$: унитарное сопряжение
        $U\Gam U^\dagger$ тождественно сохраняет все собственные значения, так
        что $\dot S_{\mathrm{unitary}}=0$ точно. Он не экспортирует логической
        энтропии --- это та же граница, что действует для всякого логически
        обратимого вычисления~\cite{bennett1973}.
  \item \textbf{Диссипатор} $\Diss$: линдбладов диссипатор Фано
        $\gamma(\mathcal{P}_{\mathrm{Fano}}(\Gam)-\Gam)$ (и опциональная
        деполяризация) \emph{повышает} $S_{\mathrm{vN}}$ со скоростью
        $\dot S_{\Diss}>0$ бит/время. Удержание когерентного состояния против
        этой утечки требует ре-экспорта
        $\Delta H_{\Diss}^{(+)}=\dot S_{\Diss}\Delta\tau$ бит за тик --- члена
        поддержания.
  \item \textbf{Регенерация} $\Regen$: движение $\Gam$ к заякоренной самомодели
        $\varphi_J(\Gam)$ \emph{понижает} $S_{\mathrm{vN}}$ на
        $\Delta H_{\Regen}$ бит за тик (создавая порядок), что требует не менее
        $k_B T\ln 2\cdot\Delta H_{\Regen}$ свободной энергии.
\end{itemize}
Суммирование даёт $E_{\mathrm{tick}}\ge k_B T\ln 2\,(\Delta H_{\Regen}+\Delta H_{\Diss}^{(+)})$,
Ур.~\eqref{eq:energy-law}.

\paragraph{Точные числа ($T=300$~К) и честный охват.}
При $T=300$~К имеем $k_B T\ln 2 = 1{,}380649\!\times\!10^{-23}\cdot 300\cdot\ln 2
= 2{,}87\times 10^{-21}$~Дж ($0{,}018$~эВ на бит). Для $\Gam\in\Dens(\mathbb{C}^7)$
точный максимум энтропии фон Неймана достигается при $\Gam=\mathbb{1}/7$:
\[
S_{\max} \;=\; \log_2 7 \;\approx\; 2{,}807355\;\text{бита}
\quad\Longrightarrow\quad
E_{\max,\text{state}} \;=\; k_B T\ln 7 \;=\; (\log_2 7)\,k_B T\ln 2
\;\approx\; 8{,}060\times 10^{-21}\;\text{Дж}.
\]
Критически важно, что $S_{\max}=\log_2 7\approx 2{,}807$~бита есть точный
потолок энтропии фон Неймана $7$-мерной матрицы плотности --- а \emph{не}
$48$~бит (вещественная размерность многообразия $\Dens(\mathbb{C}^7)$) и
\emph{не} $6{,}272$~бита (цифровое хранение $49$ элементов \texttt{complex128}
в IEEE-754). Все расчёты энтропии в битах и независимые от шага скорости
диссипации проверяются непосредственно в
\texttt{architecture/talos\_v0\_emulator.py}. Мы \emph{не} делим физические
джоули КМОП DRAM на теоретический пол Ландауэра: пределы Ландауэра относятся к
любой логически обратимой машине, тогда как мощность КМОП определяется
переключающей ёмкостью и утечками. На цифровом железе непосредственный
энергетический выигрыш \TALOS{} состоит в удержании $784$-байтного состояния в
L1/регистрах без внекристальных выборок DRAM.
